\SourceLeaf{416}%
straight line one must also specify the sense in which it is considered to be traced. Without such a distinction the direction of a straight line would always correspond to two opposite radii.

\Needspace{5\baselineskip}
\begin{center}2.\end{center}
\noindent\hspace*{1.2em}In the auxiliary circle one takes any radius as the first, or its endpoint on the circumference as the initial point, and decides in which sense arcs from it are to be counted positively (whether from left to right or the reverse). In the opposite sense the arcs are therefore regarded as negative. Thus every direction of a straight line is expressed in degrees etc. or by a number giving them in parts of the radius.

Such indications differing by $360^{0}$ or a multiple of it in fact refer to the same directions and may generally be regarded as interchangeable. In cases, however, where the manner in which a varying inclination arises is also taken into account, it may be necessary to distinguish such indications carefully.

If, for example, one has decided to count the arcs from left to right, and two straight lines $l,l'$ have corresponding directions $L,L'$, then $L'-L$ is the inclination of those straight lines. It is easy to see that insofar as $L'-L$ falls between $-180^{0}$ and $+180^{0}$, the positive or negative value also shows whether, as seen from the point of intersection, $l'$ lies to the right or left of $l$. In general this is decided by the sign of $\sin(L'-L)$.

Let $aa'$ be a portion of a curved line, and let the tangents at $a,a'$ correspond respectively to the directions $\alpha,\alpha'$. These letters are to denote the corresponding points of the auxiliary circle, and their distances along the arc from the initial point are to be denoted by $A,A'$. The magnitude of the circular arc $\alpha\alpha'$, or $A'-A$, is called the \textls[100]{Amplitudo} of $aa'$.

Comparing the Amplitudo of the arc $aa'$ with its length will give us the concept of curvature. Let $l$ be a variable point of the\space%
